\documentclass[10pt,a4paper]{amsart}
\usepackage[margin=19mm]{geometry}
\usepackage{amsmath,amssymb,mathtools}
\usepackage[scr=rsfs,cal=euler]{mathalfa}
\usepackage{xfrac}
\usepackage[unicode,hidelinks,bookmarks=false]{hyperref}
\newcommand{\ie}{i.e.\ }
\newcommand{\cf}{cf.\ }
\newcommand{\resp}{respectively\ }
\newcommand{\Subsection}{\subsection}
\providecommand{\nobreakdash}{\nobreak\hbox{-}}
\setlength{\emergencystretch}{2em}
\multlinegap0pt
\usepackage{bm}
\usepackage{bbm}
\usepackage{dsfont}
\usepackage{xspace}
\usepackage{cancel}
\usepackage{mathtools}
\usepackage{amsthm}
\makeatletter
\@ifundefined{Theorem}{\newtheorem{Theorem}{Theorem}}{}
\@ifundefined{Lemma}{\newtheorem{Lemma}[Theorem]{Lemma}}{}
\makeatother
\makeatletter
\@ifundefined{Proof}{\newenvironment{Proof}[1][Proof]{\noindent {\bf Proof.\;}}{\qed}}{}
\makeatother
\title[Some structural properties of mixed orthogonal arrays and th]{Some structural properties of mixed orthogonal arrays and their irredundancy}
\author{Maryam Bajalan, Peter Boyvalenkov, Ferruh \"Ozbudak}
\date{Source: arXiv:2603.18568v1 (CC BY 4.0)}
\begin{document}
\maketitle

\noindent\emph{Source and licence.} Some structural properties of mixed orthogonal arrays and their irredundancy. Version arXiv:2603.18568v1 (CC BY 4.0), distributed under the Creative Commons Attribution 4.0 International licence, as recorded in the arXiv licence metadata for this exact version and shown on the abstract page https://arxiv.org/abs/2603.18568v1. Full rights notice: licenses/arxiv-2603.18568-CC-BY-4.0.txt.\par\smallskip

\noindent\textit{Editorial scope.} Counted unit: the Theorem whose source label is \texttt{ma20} in the article (source0.tex lines 653--655, source label \texttt{ma20}), delivered together with the article's own complete original proof of it (source0.tex lines 656--679, one \texttt{proof} environment of 24 lines). No mathematics is written, repaired or re-derived: statement and proof are the article's own text line for line. Required context retained uncounted: parity (lines 224--228); ma14 (lines 238--244); surjective (lines 596--603); coset (lines 634--637). Every definition, display and label of those lines is retained unchanged. Omitted from this package and not redistributed: 1--223, 229--237, 245--595, 604--633, 638--652, 680--1044. Nothing inside a retained excerpt is omitted or altered: every retained body is delivered exactly as the article's lines are written, so no omission receipt of any kind accompanies this package. Citations: the retained excerpts carry no citation command at all, so no bibliography entry of the article is transcribed and no reference is redistributed. Reference adaptation: none was needed and none was made; every reference inside the delivered unit resolves inside it, so no unresolved reference or citation is produced. Printed slips kept literal: no printed word, symbol, hypothesis or number of the counted statement or of its proof was altered. Layout and class: the article's own class is \texttt{article} with packages including natbib, fontenc, graphicx, inputenc, amsmath, mathrsfs, mathtools, amsthm, amsfonts, amssymb, epsfig, color, epstopdf, yfonts; the wrapper is the lane's pinned \texttt{amsart} editorial wrapper at 10pt on A4, which restates the theorem-like environments the retained text needs and adds the article's own macro definitions that the retained text needs (none), with extra wrapper packages (bm, bbm, dsfont, xspace, cancel, mathtools). The delivered numbering is the unit's own. Assets: no retained span contains a figure environment, a graphics-inclusion command, a PDF-page inclusion, a table or a TikZ picture; no image file is copied into this package, the package declares no asset dependency and no image file is redistributed with it. Licence: version arXiv:2603.18568v1 of this article is distributed under the Creative Commons Attribution 4.0 International licence, as recorded in the arXiv licence metadata for this exact version and shown on the abstract page https://arxiv.org/abs/2603.18568v1. A full rights notice is delivered at licenses/arxiv-2603.18568-CC-BY-4.0.txt. Characters: every retained excerpt is pure ASCII, so the delivered engine, XeLaTeX with its default Latin Modern text font, reports no missing character.
\begin{align}\label{parity}
    H = \begin{bmatrix}
H_1 \mid H_2 \mid \dots \mid H_s
\end{bmatrix},
\end{align}

\begin{Lemma}\label{ma14}
Let $C\subseteq\mathbb F_q^n$ be an $\mathbb F_q$-linear error-block code of type $\pi=[n_1]\cdots[n_s]$ with the minimum $\pi$-distance $d_{\pi}$ and with the parity-check matrix $H$ given in the form \eqref{parity}.
\begin{enumerate}
    \item If, for every choice of $t$ blocks of $H$, the union of their columns is an $\mathbb F_q$-linearly independent set, then $d_{\pi}(C)\ge t+1.$
    \item If there exists a set of $t+1$ blocks of $H$ whose union of columns is $\mathbb F_q$-linearly dependent, then $d_\pi(C)\le t+1.$
\end{enumerate}
\end{Lemma}

\begin{Lemma}\label{surjective}
Let $C$ be an $\mathbb F_q$-linear error-block code in $ \mathbb{F}_q^n =\mathbb{F}_q^{n_1}\times \mathbb{F}_q^{n_2} \times \ldots\times \mathbb{F}_q^{n_s}$ with the Euclidean dual $C^\perp$. Let $\{i_1,\ldots,i_t\}\subseteq\{1,\ldots,s\}$ be a fixed selection of $t$ blocks, and consider the $\mathbb F_q$-linear projection map
        \begin{align*}
         \mu : C &\to \mathbb{F}_{q}^{n_{i_1}} \times \cdots \times \mathbb{F}_{q}^{n_{i_t}}\\  
          \mathbf c=(\mathbf c_1,\ldots,\mathbf c_s)&\mapsto(\mathbf c_{i_1},\ldots,\mathbf c_{i_t}).
        \end{align*}
If $\mu$ is not surjective, then $d_\pi(C^\perp)\le t$.
\end{Lemma}

\begin{Lemma} \label{coset}
Let $A$ and $B$ be $\mathbb F_q$-vector spaces, and let $f: A \to B$ be a surjective $\mathbb F_q$-linear map. Then for every $b \in B$,
$$|f^{-1}(b)| = |\ker f|.$$
\end{Lemma}

\begin{Theorem}\label{ma20}
Let $C$ be an $\mathbb F_q$-linear $M \times s$ array whose rows lie in $\mathbb{F}_{q^{n_1}}\times\mathbb{F}_{q^{n_2}}\times\cdots\times\mathbb{F}_{q^{n_s}}$. If every set of $t$ columns of $C$ is $\mathbb F_q$-linearly independent, then $C$ is an $\mathbb F_q$-linear mixed orthogonal array $\mathrm{MOA}\left(M, s, (q^{n_1}, \dots, q^{n_s}), t\right)$.
\end{Theorem}

\begin{proof}
Since $\rho$ is an $\mathbb F_q$-isomorphism, $t$ columns $V_1, V_2, \ldots, V_t$ of $C$ are $\mathbb F_q$-linearly independent if and only if the union of the columns of the blocks $\rho(V_1), \rho(V_2), \ldots, \rho(V_t)$ is an $\mathbb F_q$-linearly independent set. Therefore, by Lemma~\ref{ma14}, we obtain $d_\pi((\rho(C))^\perp)\geq t+1$.\\
For a fixed selection of $t$ columns $\{i_1,\ldots,i_t\}$ of $C$, we define the projection maps $\mu'$ and $\mu$ as shown in the following diagram:
  \begin{align}
\begin{array}{ccc}
C & \xrightarrow{\quad \mu'\quad} & \mathbb{F}_{q^{n_{i_1}}}\times\mathbb{F}_{q^{n_{i_2}}}\times\cdots\times\mathbb{F}_{q^{n_{i_t}}} \\
\Big\downarrow{\scriptstyle \rho} &  & \Big\downarrow{\scriptstyle \rho} \\
\rho(C) & \xrightarrow{\quad \mu \quad} & \mathbb{F}_{q}^{n_{i_1}}\times\mathbb{F}_{q}^{n_{i_2}}\times\cdots\times\mathbb{F}_{q}^{n_{i_t}}.
\end{array}
\end{align} 
Since $d_\pi((\rho(C))^\perp)\ge t+1$, Lemma~\ref{surjective} implies that $\mu$ is surjective. Since $\rho$ is an isomorphism, it follows that $\mu'$ is also surjective, and hence $ |\mu'(C)|=q^{n_{i_1} + \cdots + n_{i_t}}.$ By the first isomorphism theorem,
$$ q^{n_{i_1} + \cdots + n_{i_t}}=|\mu'(C)|= \frac{|C|}{|\ker(\mu')|},$$
and hence
\begin{align}\label{First11}
    |\ker(\mu')| = \frac{M}{q^{n_{i_1} + \cdots + n_{i_t}}}
\end{align}
Using Lemma \ref{coset} and equation \eqref{First11}, for any  $\boldsymbol{c}' \in  \mathbb{F}_{q^{n_{i_1}}} \times \mathbb{F}_{q^{n_{i_2}}}\times \cdots \times \mathbb{F}_{q^{n_{i_t}}}$,  we have 
$$
|(\mu')^{-1}(\boldsymbol{c}')|=\frac{M}{q^{n_{i_1} + \cdots + n_{i_t}}}.
$$
Hence there are $\frac{M}{q^{n_{i_1} + \cdots + n_{i_t}}}$ codewords in $C$ that map to the  $t$-tuple $\mathbf{c}'$ under the map $\mu'$. \\
Therefore, for every selection $\{i_1,\ldots,i_t\}$ of $t$ columns of $C$, each $t$-tuple appears exactly
$$\frac{M}{q^{n_{i_1} + n_{i_2} + \cdots + n_{i_t}}}.$$  Therefore, $C$ is an $ \mathrm{MOA}\left(M, s, (q^{n_1},\dots, q^{n_s}), t\right)$.
\end{proof}

\end{document}
